\newpage
\chapter*{参考文献}
\addcontentsline{toc}{chapter}{参考文献}
只列出作者直接阅读过且在正文中被引用过的文献资料，本专业教科书一般不能作为参考文献，除个别专业外，一般应有外文参考文献。参考文献一律放在论文参考文献部分中，不应放在章节末尾或页面脚注中。参考文献一般不应少于15篇。“参考文献”四字采用不编号章标题样式。内容样式与正文相同。
根据GB/T 7714—2015的要求书写参考文献，可采用顺序编码制，也可采用著者-出版年制，但全文必须统一，各院系内部需统一。顺序编码制的方括号和编号采用宋体。
参考文献作者的著录方法如下：

（1）3人以下全部著录，3人以上可只著录前3人，后加“，等”，外文用“, et al ”“ et al ”不必用斜体；

（2）责任者之间用“，”分隔；

（3）责任者姓名一律采用姓前名后的著录形式。

参考文献中的标点符号：中文文献采用中文、全角、中文标点输入法输入（下角点“.”除外），标点后接排后续内容；英文文献采用英文、半角、英文标点输入法输入，标点后空一格编排后续内容。

\printbibliography[heading=none]